Beyond benchmark-style discussion, our project now contains a concrete application pipeline built around historical airplane crash records. This case study is valuable because it forces the method to operate on messy real data rather than on idealized toy examples. The notebook \texttt{airplane\_crashes\_gmm\_imputation.ipynb} uses the helper module \texttt{src/airplane\_crashes\_pipeline.py} to download the dataset, engineer numeric features, inject MCAR missingness, impute the missing entries with our GMM routine, and generate visual artifacts that can be reused directly in the paper.

\subsection{Pipeline Design}
The application pipeline transforms the raw crash records into six numeric variables:
\[
    \texttt{Aboard},\ \texttt{Fatalities},\ \texttt{Ground},\ \texttt{Year},\ \texttt{Month},\ \texttt{Day}.
\]
Rows with all-numeric features missing are discarded, while fully observed rows are retained as the reference matrix used for controlled evaluation. Missingness is then injected synthetically at rate $20\%$ so that the imputation quality can be assessed objectively on the masked entries.

For the configuration currently stored in the notebook outputs, the pipeline reports:
\begin{itemize}
    \item $5236$ usable rows after preprocessing;
    \item $6$ numeric features in the modeling table;
    \item GMM imputation with $\scal{K}=4$ mixture components;
    \item RMSE $= 20.0588$ on the masked entries.
\end{itemize}

This section matters because it upgrades the paper from ``we implemented the algorithm'' to ``we deployed it on a reproducible domain-specific workflow with measurable reconstruction performance.''

\subsection{Artifacts for the Paper}
Three visual artifacts generated by the notebook are especially valuable for the report because they illustrate complementary aspects of the workflow.

\begin{figure}[t]
    \centering
    \includegraphics[width=\linewidth]{applications/missingness_ratio.png}
    \caption{Feature-wise missingness profile for the airplane-crash feature table before synthetic corruption. This plot supports the data-quality narrative and motivates explicit imputation.}
    \label{fig:aircraft-missingness}
\end{figure}

\begin{figure*}[t]
    \centering
    \includegraphics[width=0.92\textwidth]{applications/distribution_comparison.png}
    \caption{Distribution comparison between ground-truth numeric features and GMM-imputed features in the airplane-crash application. The overlap is not perfect, but it shows that the imputation preserves the broad marginal structure of the data.}
    \label{fig:aircraft-distribution}
\end{figure*}

\begin{figure}[t]
    \centering
    \includegraphics[width=\linewidth]{applications/pca_clusters.png}
    \caption{PCA projection of the imputed airplane-crash feature matrix colored by the learned mixture assignments. The plot is useful not because the clusters are guaranteed to be semantically pure, but because it reveals the geometry that the model imposes after imputation.}
    \label{fig:aircraft-pca}
\end{figure}

\subsection{Why This Application Strengthens the Paper}
This application is not just decorative. It contributes to the credibility of the report in three concrete ways:
\begin{itemize}
    \item it demonstrates that the method is usable on a real tabular dataset with heterogeneous scales and event-based features;
    \item it produces a quantitative imputation score, namely RMSE on masked entries, rather than relying only on clustering metrics;
    \item it yields visual evidence that can be inspected independently of the numerical tables.
\end{itemize}

In other words, the airplane-crash pipeline helps the paper move from method exposition to method deployment.
